\label{sec:discussion}

The central finding of this paper is that \textsc{bisect} algorithmic maps produce
incumbency outcomes statistically indistinguishable from the incumbency-neutral baseline,
while enacted maps produce outcomes systematically below that baseline on pairing rate
and above it on safe-seat fraction. This finding is analytically expected but empirically
important: it demonstrates that the discrepancy between algorithmic and enacted maps is
not an artefact of geographic constraints (which affect both maps equally) but of
deliberate mapmaker choices about where to draw district boundaries relative to incumbent
addresses and partisan voter concentrations.

\subsection{The Neutrality Argument}

We call this the \textit{neutrality argument}: the \textsc{bisect} algorithm is
incumbency-neutral not in the sense that it produces no incumbency consequences, but in
the sense that the incumbency consequences it produces are a predictable function of
geography alone. Given the geographic distribution of a state's voters and the location
of incumbent homes within that geography, the algorithm produces a specific incumbency
outcome, and that outcome is well characterised by the incumbency-neutral random baseline.
No strategic choice about incumbency protection is embedded in the algorithm; the only
embedded choices are about population balance, contiguity, and compactness.

This distinction matters legally. When a state legislator testifies that a proposed
algorithmic map ``unfairly pairs'' two incumbents, the appropriate response is not that
pairing was intended, but that pairing was the inevitable result of drawing compact,
population-balanced districts in a geography where those two incumbents live geographically
close to each other. The enacted map's zero pairing rate, by contrast, is not a function
of geography---it is a function of deliberate boundary choices that required access to
incumbent address data and the political will to use that data to protect incumbents.

\subsection{Incumbency Outcomes as Purely Geometric}

The three mechanisms identified in Section~\ref{sec:introduction} all reduce to the
same geometric observation: \textsc{bisect} draws district boundaries by optimising a
cost function that has nothing to do with incumbent locations, and the resulting
boundaries happen to intersect the geographic distribution of incumbents in ways that
produce pairing, partisan composition change, and constituent reassignment. These
outcomes are fully determined by two inputs: the algorithm's geographic optimisation
and the geographic distribution of incumbent homes. Neither input involves any strategic
assessment of incumbency; the outcomes are therefore purely geometric, not strategic.

Enacted maps, by contrast, produce incumbency outcomes that cannot be explained by
geography alone. The enacted maps' zero pairing rates and high safe-seat fractions require
an explanation that goes beyond compact geography---they require the hypothesis that
mapmakers actively used incumbent addresses and partisan data to engineer protective
boundaries. I.1 and I.2 provide formal statistical tests of this hypothesis; both
papers reject the null that enacted pairing rates and safe-seat fractions are consistent
with random redistricting.

\subsection{Implications for Expert Witness Testimony}

The neutrality argument has direct implications for redistricting litigation. When an
expert witness proposes \textsc{bisect} maps as a neutral baseline in a gerrymandering
case, opposing counsel may argue that the algorithm ``unfairly'' disrupts incumbency by
producing more pairings and more open seats than the enacted map. The neutrality argument
provides the response: the algorithmic map produces incumbency outcomes consistent with
what a neutral random process would produce given the state's geography. The enacted map's
lower pairing rate and higher safe-seat fraction are not geographic necessities; they are
choices. Any map that achieves a lower pairing rate than the algorithmic baseline does so
by using incumbent location data---a choice, not a necessity.

This framing also addresses the VRA tension noted in Section~\ref{sec:background}. When
a \textsc{bisect} map pairs a minority incumbent with an Anglo incumbent, it does so
because their homes happen to fall on the same side of a geographically optimal boundary,
not because the algorithm targeted either incumbent. The pairing is a geographic accident;
the remedy, if any, lies in post hoc adjustment of the boundary to honour VRA requirements,
not in attributing strategic intent to an algorithm that had none.
